% 第 13 题(2) 解析：圆形磁场区与正方形设计区（按原图矢量重绘）
\begin{tikzpicture}[font=\small, >=Stealth, scale=0.92]
  % 坐标系
  \draw[->] (-1.3,0) -- (5.4,0) node[right] {$x$};
  \draw[->] (0,-3.4) -- (0,3.5) node[above] {$y$};
  \node[below left] at (0,0) {$O$};
  % 上方圆形磁场区（圆心 A(0,r1)，半径 r1）
  \draw[thick] (0,1.6) circle (1.6);
  \fill (0,1.6) circle (1.5pt) node[right,font=\scriptsize] {$A(0,r_{1})$};
  % 入射粒子流（宽度 2r1，沿 +x）
  \draw[->,thick] (-2.4,3.2) -- (-0.05,3.2);
  \draw[->,thick] (-2.4,2.0) -- (-0.05,2.0);
  \draw[->,thick] (-2.4,0.8) -- (-0.05,0.8);
  \draw[->,thick] (-2.4,0.0) -- (-1.3,0.0);
  \draw[<->] (-2.2,0.0) -- (-2.2,3.2);
  \node[left,font=\scriptsize] at (-2.2,1.6) {$2r_{1}$};
  % 下方正方形虚线框（中心 C(0,-r2)，边长 2r2）
  \draw[thick,dashed] (-2.0,-3.2) rectangle (2.0,0.8);
  \fill (0,-1.2) circle (1.5pt) node[right,font=\scriptsize] {$C(0,-r_{2})$};
  % 下方磁场区（圆，半径 r2，画在正方形内）
  \draw[thick] (0,-1.2) circle (2.0);
  % 磁场符号
  \foreach \x in {-1.4,-0.7,0,0.7,1.4} {
    \foreach \y in {-2.9,-2.2,-1.5,-0.8,-0.1} {
      \draw (\x,\y) ++(-0.06,-0.06) -- ++(0.12,0.12);
      \draw (\x,\y) ++(-0.06,0.06) -- ++(0.12,-0.12);
    }
  }
  % 出射粒子流（宽度 2r2）
  \draw[->,thick] (2.05,-3.2) -- (4.6,-3.2);
  \draw[->,thick] (2.05,0.8) -- (4.6,0.8);
  \draw[<->] (4.2,-3.2) -- (4.2,0.8);
  \node[right,font=\scriptsize] at (4.25,-1.2) {$2r_{2}$};
\end{tikzpicture}
